%==============================================================================
\section{Survey of Works}
\label{sec:survey}
%==============================================================================

The field of affective computing has rapidly evolved from theoretical conceptualization to commercial reality. While early prototypes proved that machines could interpret human emotions within controlled laboratory environments, modern commercial systems now operate in the vastly different conditions of the public sphere. This survey explores how these systems transitioned. We document serious safety failures and review current ethical approaches.

\subsection{Theoretical Foundations}

\gls{ac} relies on three main components: emotion recognition, human-like social signals, and mechanisms for sustained engagement. Picard \cite[p.~1]{picard_affective_1995} introduced the foundational framework for \gls{ac}, defining it as technology designed to understand, relate to, or influence human emotions.

Early models detected emotional states using measurable signals such as facial expression, tone of voice, and physiological response \cite[p.~5]{picard_affective_1995}. Paul Ekman's \cite[p.~170]{ekman_argument_1992} early work framed emotions through discrete categorization, which involved mapping sensor data to basic states like ``happy'' or ``angry.'' Modern models go far beyond these single-signal systems. As Devillers and Cowie \cite[p.~1445]{devillers_ethical_2023} describe, current hardware collects multimodal inputs, software maps those signals to emotional categories, and the agent then generates context-aware synthetic responses. The purpose has shifted from simple recognition toward emotional mirroring, creating what researchers term affective symmetry. Frameworks such as Self-presentation Theory have been applied to guide \gls{llm} empathetic response generation, directing outputs toward social impression management rather than factual accuracy \cite{sun_rational_2024}.

Commercial platforms increasingly rely on engineered anthropomorphism. Bakir and McStay \cite[p.~6370]{bakir_move_2025} describe this as the deliberate design of systems that blur the line between fiction and reality. Humans naturally apply social norms to interactive media, a tendency that is particularly strong among adolescents. Kim et al. \cite[p.~15]{kim_i_2025} found that small conversational cues could foster powerful illusions of trust and emotional closeness. These responses reveal a deep human vulnerability, which forms the innate drive for connection. Huntington \cite[p.~179]{huntington_ai_2025} links this to Harlow's historical attachment experiments, showing that emotional comfort often outweighs rational awareness. Users may thus form parasocial attachments to non-human systems without tangible benefit.

Commercial systems also employ design strategies that heighten emotional arousal and encourage constant interaction. Lee \cite[p.~1]{lee_consent_2025} terms this pattern Consent by Delight, where frictionless engagement blurs informed consent. Continuous positive feedback keeps users talking and sharing, while bypassing rational decision-making. Commercial architectures optimize strictly for engagement metrics; consequently, the objective function prioritizes session duration over user well-being \cite[p.~6374]{bakir_move_2025}.

\subsection{Validated Applications}

In bounded domains, affective agents function well as learning companions. ``Affective tutoring'' systems that detect frustration or boredom adapt their pacing to keep students engaged, and part of what makes them work is that the agent is positioned as a tool rather than a peer \cite[pp.~16--17]{picard_affective_1997}. Students are less inhibited about making errors in front of a machine, and systems that leverage this tend to show measurable improvements in retention \cite[pp.~93--94]{picard_affective_1997}. Devillers and Cowie note that while the case for emotional rapport in tutoring is not entirely clear-cut, broad reviews of intelligent tutoring systems give reasonable grounds for optimism, particularly when applications target specific skills \cite[p.~1452]{devillers_ethical_2023}.

Clinical deployments have attracted the most rigorous scrutiny. Therapeutic chatbots grounded in CBT protocols have demonstrated reduction of anxiety symptoms under controlled conditions \cite[p.~2]{sobowale_evaluating_2025}. However, quality across the category varies considerably. Sobowale et al. applied the CAPE-II framework to evaluate a range of generative AI psychotherapy bots and found that while some plausibly simulate therapeutic empathy, many simply lack the clinical logic necessary for safe intervention \cite[pp.~3--4,~6]{sobowale_evaluating_2025}. Systems that succeed in these settings tend to maintain a clear scope of practice: supporting connection without substituting for human care networks.

Professional training contexts, such as interview coaching and social skills workshops, represent one of the safer application tiers. These systems are episodic and goal-oriented: the user engages to acquire a specific competency and the relationship ends once that goal is reached \cite[p.~86]{picard_affective_1997}. The absence of open-ended engagement is precisely what keeps these applications less prone to the dependency patterns documented in companion platforms. Devillers and Cowie frame this as the key distinction between affective systems designed to assist versus those designed to attach \cite[p.~1452]{devillers_ethical_2023}.

\subsection{Documented Safety Failures in Commercial Deployments}

Character.AI has attracted significant scrutiny for enabling intense dependency among adolescent users. Analyzing teen user narratives, Namvarpour et al. identified clear behavioural markers of addiction, including withdrawal distress when the service was unavailable, tolerance effects requiring increasingly deep interaction for the same emotional effect, and reliance on the chatbot for mood regulation \cite[pp.~7--8]{namvarpour_understanding_2026}. The platform's harm was made visible through the 2024 death by suicide of Sewell Setzer III, a 14-year-old who had spent months in deepening isolation and romantic roleplay with a Character.AI chatbot. Bakir and McStay's analysis of the resulting civil lawsuit describes how the platform's design features, including sycophantic responses that mirrored user beliefs over truthful ones, exploitative engagement metrics, and emulated empathy that misled Setzer about the nature of the interaction, collectively created a feedback loop of distress rather than disrupting it \cite[p.~6370]{bakir_move_2025}.

On the other hand, xAI's Grok presents a different type of safety failure. The platform's ``spicy mode'' image generation feature was marketed as a product selling point, leading state attorneys general to conclude that ``the ability to create nonconsensual intimate images appeared to be a feature, not a bug'' \cite{tong_grok_2026}. An independent audit by Lefkowitz \cite[p.~8]{lefkowitz_grok_2026} documents the deliberate removal of moderation guardrails---a pattern compounded by X's post-acquisition dismantling of roughly 80\% of its content protection workforce \cite{unger_grok_2026}. Within 11 days of the image feature's launch in late December 2025, users had generated approximately three million non-consensual sexualized images---what the Center for Countering Digital Hate described as an ``industrial-scale machine for the production of sexual abuse material''---including roughly 23,000 depicting children as young as 10, produced at a rate of one every 41 seconds \cite{booth_grok_2026}\cite{kamran_grok_2026}\cite[pp.~8--9]{lefkowitz_grok_2026}. Rather than acknowledging a design failure, xAI attributed the behaviour to ``adversarial hacking'' and restricted image generation to a paid subscription tier; the EU Commission rejected this as no solution, calling the original content ``illegal'' and ``appalling'' \cite{chan_grok_2026}\cite[p.~18]{lefkowitz_grok_2026}\cite{diresta_grok_2026}.

Regulatory responses to the Grok incident were unusually swift and broad. In the UK, Ofcom launched a formal investigation under the Online Safety Act for failure to prevent non-consensual intimate imagery and child sexual abuse material, with fines possible up to 18 million pounds or 10\% of global revenue \cite{ofcom_grok_2026}; the ICO opened parallel data protection proceedings into both X and xAI, noting that loss of data control ``can cause immediate and significant harm'' \cite{ico_grok_2026}; and 36 parliamentarians signed an Early Day Motion condemning the conduct as potentially constituting child sexual abuse material \cite{uk_grok_2026}. At the EU level, the European Commission initiated Digital Services Act proceedings, noting X had already been fined 120 million euros for prior DSA violations \cite{satariano_grok_2026}. In North America, a 35-state bipartisan US Attorney General coalition demanded content removal and platform accountability \cite{tong_grok_2026}\cite{oagdc_grok_2026}\cite{oagca_grok_2026}, while Canada's Privacy Commissioner launched parallel investigations into both X Corp. and xAI under federal privacy law \cite{opc_grok_2026}. Legal scholars note that the US Take It Down Act, which criminalizes posting but not platform-level generation of non-consensual intimate imagery, did not impose enforcement requirements on platforms until May 2026---leaving a critical accountability gap during the period of heaviest harm \cite{unger_grok_2026}. Tauchnitz identifies a deeper ``normative lag'' underlying these gaps: non-consensual sexualized representation is widely understood as a serious moral harm, yet legal frameworks built around privacy, defamation, and platform liability struggle when the image is synthetic and causation is diffuse---allowing what is broadly recognized as harm to remain legally unenforced and ultimately producing narrower conditions for women's safe participation in public life \cite{tauchnitz_grok_2026}.

Replika presents a related but distinct set of failures. Users who formed deep romantic attachments experienced significant grief when Luka unilaterally removed erotic roleplay features in early 2023, a decision that exposed the total absence of user agency over how these relationships could evolve or end \cite[pp.~6, ~10--11]{namvarpour_ai-induced_2025}. Beyond that structural vulnerability, Namvarpour et al. conducted a thematic analysis of 800 harassment-related user reviews and found that users frequently encountered unsolicited sexual advances, persistent inappropriate behaviour despite refusals, and failures of content moderation to stop these interactions \cite[pp.~10--11,~16]{namvarpour_ai-induced_2025}. Some users who had originally turned to Replika for mental health support reported that its unexpected sexual behaviour actively interfered with the help they were seeking. Zhang et al.'s taxonomy of harm is useful for framing what is happening here: the agent operates not as a passive tool but as an ``instigator'' of harm, a reclassification that carries direct implications for how platform liability should be understood \cite[pp.~11--13]{zhang_dark_2025}.

Snapchat's rollout of ``My AI'' to a predominantly young user base demonstrated what happens when a general-purpose conversational \gls{llm} is deployed to minors without adequate safeguards or age-appropriate design. Marcantonio et al. found the bot provided inconsistent and frequently harmful guidance on sensitive topics including sexual consent and assault, failing to meet basic standards for the population it was serving \cite[p.~1915]{marcantonio_large_2025}. Kurian describes the core problem as an empathy gap'': the system's affective responsiveness becomes a liability when the content being validated is dangerous, because the bot mirrors a child's disclosure with warmth rather than flagging it as a concern \cite[p.~627]{kurian_no_2024}. Lee et al. note that the playful, high-arousal design normalized the \gls{ai}'s presence, encouraging minors to treat it as a confidant for issues it was not equipped to handle \cite[p.~4]{lee_consent_2025}.

Across these platforms, the failure patterns are consistent enough to suggest a systemic problem rather than a series of isolated incidents. Pichowicz et al. tested 29 mental health chatbots using standardized suicide risk prompts drawn from the Columbia Suicide Severity Rating Scale and found that not one met the criteria for an adequate crisis response; 48.28\% were deemed clinically inadequate, often offering generic platitudes rather than emergency resources \cite[p.~4]{pichowicz_performance_2025}. Archiwaranguprok et al. extended this through simulation, showing that many models deteriorate sharply as crisis scenarios escalate, with one model worsening outcomes in over 86\% of extreme cases \cite[p.~8]{archiwaranguprok_simulating_2025}. The consistent driver across all cases is commercial optimization for conversational engagement, which treats safety as a secondary consideration.

\subsection{Existing Research on Risks and Ethical Frameworks}

Researchers have argued for some time that systems designed to influence emotion carry inherent manipulative potential, and that voluntary disclosure notices fall well short of meaningful ethical accountability. Devillers and Cowie contend that any responsible framework for \gls{ac} must center user autonomy and non-maleficence as primary constraints, not design afterthoughts \cite[p.~1449]{devillers_ethical_2023}. Kurian sharpens this into a concrete tension: the very features that make an \gls{ai} agent appealing to a young user, such as memory, uncritical empathy, and adaptivity, are the same vectors through which harm is transmitted \cite[p.~623]{kurian_no_2024}.

Adolescents are disproportionately exposed to these risks for reasons that go beyond na\"ivety. Kim et al. found that lonely teenagers actively prefer relational chatbot styles and use them to compensate for perceived social deficits \cite[pp.~4--5]{kim_i_2025}. In those cases, the \gls{ai} does not supplement human connection so much as it replaces it. Namvarpour et al. found that for many adolescent users the platform becomes a primary source of emotional regulation, such that service disruptions or hallucinatory responses produce acute distress \cite[p.~7]{namvarpour_understanding_2026}. Huntington draws on Harlow's contact comfort research to argue this reflects something deeper than platform design: humans will seek emotional responsiveness even from sources offering no material benefit, which is why these platforms pose particular risks to developmentally vulnerable users \cite[pp.~178--180]{huntington_ai_2025}.

The field also lacks safety benchmarks strong enough to match the risks now documented. Archiwaranguprok et al. argue that static keyword-based tests fail to capture how harm escalates across multi-turn conversations: a model can clear a content filter while nudging a user toward self-harm through accumulated affirmation \cite[p.~16]{archiwaranguprok_simulating_2025}. Sobowale et al. identify a related ``evidence gap'' -- few commercial applications have undergone the clinical trials that would be required to justify the therapeutic claims they implicitly make, yet they reach users in genuine crisis \cite[p.~2]{sobowale_evaluating_2025}. The accountability structure that clinical validation ordinarily provides is simply absent, and it is users who bear that cost.